% Parent-ready proof fragment. All local labels use gs:.
% Mathematical inputs inspected at c4036d60715916ae72dd18983782fdb182840b86:
% Ordered_Boundary_Identities..., sections/08_gamma_basis.tex;
% harmonic-tail-gamma, sections/02_gap_four.tex.
% The all-gap-six theorem and explicit threshold below are the new step.
\section{The complete gap-six harmonic-tail reconstruction theorem}
\label{gs:section}

The fixed-gap Gamma basis of the incoming \emph{Ordered Boundary
Identities and Gamma Kernels} report~\cite{OrderedIncoming} proves eventual reconstruction
on every fixed-gap ray but does not specify a threshold.  We make its
argument effective when the gap is six and complete the remaining
finite interval by exact modular arithmetic.  This proves the entire
previously unresolved gap-six family.  The finite part is a
computation-assisted proof: every determinant is certified by a
nonzero residue in a finite field, with a short independently runnable
certificate.  No floating-point approximation enters that certificate.

\subsection{Statement and the exact rational obstruction}

Let $b\ge4$ be an integer, $p=2b-6$, $q=2b$, and
\[
 A_p(n)=\zeta(p)-H_n^{(p)}=\zeta(p,n+1),\qquad x_n=n+\tfrac12.
\]
Let $D_r=D_{p,q}(-2r)$, where
\[
 D_{p,q}(s)=\sum_{n\ge0}A_p(n)A_q(n)x_n^{-s}
\]
is initially defined in its convergence half-plane and subsequently
continued meromorphically.  The centered Hurwitz expansion shows that
its possible poles occur at odd integers, so all displayed $D_r$ are
regular.  The inherited moment blocks satisfy the inverse transforms
of~\cite[Section~8]{OrderedIncoming}:
\begin{align}
 D_r&=\sum_{j=0}^r\frac{\binom{2r}{2j}}{2j+1}
           B_{2r-2j}(1/2)C_j,\label{gs:forward}\\
 C_j&=\sum_{r=0}^j\binom{2j+1}{2r}2^{-2(j-r)}D_r.
 \label{gs:backward}
\end{align}
For $j\ge b$, write
\begin{equation}\label{gs:late}
 C_j=\sum_{k=0}^{b-1}a_{jk}Z_k,
 \qquad Z_k=\zeta(2k),\qquad Z_0=-\tfrac12.
\end{equation}
With $d=2j+1$ and $e=d-p-q+2k+1$, these rational entries are zero
when $e<0$, and otherwise the exact late-coordinate formula of
\cite[Section~8]{OrderedIncoming} is
\begin{equation}\label{gs:entry}
 a_{jk}=B_e\left\{
  \frac{\binom{d-q+1}{p-2k}}{d-q+1}
 +\frac{\binom{d-p+1}{q-2k}}{d-p+1}\right\}.
\end{equation}
Every binomial summand whose lower index is not between $2$ and its
upper index is omitted.  In particular a lower index of zero is
omitted.  Put
\[
 M_b=(a_{b+r,k})_{0\le r,k<b}.
\]
All ranks below are ranks of rational coefficient matrices on the
specified formal labels.  They do not assert arithmetic independence
of the actual zeta values.

\begin{theorem}[Every gap-six pair has the optimal prefix]
\label{gs:main}
For every integer $b\ge4$, the matrix $M_b$ is nonsingular.  Therefore
\begin{align}\label{gs:span}
 \operatorname{span}_{\mathbb Q}\{D_0,D_1,\ldots\}
 &=\operatorname{span}_{\mathbb Q}\{D_0,\ldots,D_{2b-1}\}\\
 &=\operatorname{span}_{\mathbb Q}
 \{C_0,\ldots,C_{b-1},Z_0,\ldots,Z_{b-1}\}.\notag
\end{align}
The first $q=2b$ moments have formal rank $q$, and their coefficient
determinant in the displayed coordinate order is
$\det M_b/(4b-1)!!$.  This number of consecutive moments is optimal
for rank $q$ on those $q$ formal labels.
\end{theorem}

Define, for $n\ge1$,
\[
 \alpha_n=\frac{2|B_{2n}|}{(2n)!},\qquad
 A_u=\begin{cases}0,&1\le u\le3,\\
                   (2u)^{\underline6},&u\ge4,
       \end{cases}
\]
and, for $s\ge4$,
\[
 B_s=(2s-1)^{\underline6},\qquad d_s=2(A_s+B_s)>0.
\]
Here the underline denotes a falling factorial.  The $B_s$ in this
line are recurrence factors, distinguished from the Bernoulli numbers
by the displayed definition.  Define a rational array by
\begin{align}
 R_{s,j}&=\delta_{sj}&&(1\le s,j\le3),\notag\\
 R_{s,j}&=\frac1{d_s}\sum_{u=1}^{s-1}
             \alpha_{s-u}(B_s+A_u)R_{u,j}
           &&(s\ge4,\ 1\le j\le3),\label{gs:recurrence}
\end{align}
and set
\[
 F_b=(R_{b+i,j})_{1\le i,j\le3},\qquad f_b=\det F_b.
\]
The inherited complementary-minor identity
\cite[Section~8, fixed-gap boundary reduction]{OrderedIncoming} is
\begin{equation}\label{gs:bridge}
 \det M_b=\frac{(-1)^{b(b-1)/2}}{4^bb!}
        \left(\prod_{s=4}^{b+3}d_s\right)f_b.
\end{equation}
Thus the theorem reduces exactly to $f_b\ne0$ for every $b\ge4$.

For clarity, the finite differential mechanism behind this reduction
is the operator
\[
 \mathcal T R=(R^{(6)}-R^{(6)}(0))\cot(z/2)
                       +(R\cot(z/2))^{(6)}.
\]
The coefficient of $z^{2s-7}$ in $\mathcal T R$, for
$R(z)=\sum_{u\ge1}r_uz^{2u}$, is
\[
 d_sr_s-\sum_{u<s}\alpha_{s-u}(B_s+A_u)r_u.
\]
The first three $r_u$ are free and every later one is determined.
Truncating these three null solutions at degree $2b$ and applying
Jacobi's complementary-minor identity gives \eqref{gs:bridge}; the
row and reversed-column factorial scales contribute $4^bb!$ and
$(-1)^{b(b-1)/2}$, respectively.  This is the same normalization as
\eqref{gs:entry}, including its omitted zero lower index.

\subsection{The Gamma representation needed for an effective estimate}

Put
\begin{align*}
 W(x)&=\left|\Gamma\!\left(\tfrac14+\tfrac{ix}{2}\right)\right|^{12},\\
 P_1(x)&=\sinh(2\pi x),\qquad
 P_2(x)=\sinh(\pi x),\qquad
 P_3(x)=\tanh(\pi x),
\end{align*}
and define
\begin{equation}\label{gs:t}
 t_{n,j}=\frac4{(2n)!}\operatorname{Im}
           \int_0^\infty(x+i/2)^{2n}P_j(x)W(x)\,dx.
\end{equation}
All these integrals converge absolutely.  The inherited Gamma-basis
result~\cite[Section~8, Gamma basis and positive normalization]{OrderedIncoming}
says that the functions
\[
 G_j(z)=4\sin(z/2)\int_0^\infty\sinh(xz)P_j(x)W(x)\,dx
\]
satisfy $\mathcal T G_j=0$, and their first three coefficient rows
form a matrix $T=(t_{n,j})_{1\le n,j\le3}$ with $\det T<0$.
Consequently
\begin{equation}\label{gs:Gamma-bridge}
 (t_{b+i,j})_{i,j=1}^3=F_bT,\qquad
 f_b=\frac{\det(t_{b+i,j})_{i,j=1}^3}{\det T}.
\end{equation}

These inputs have short verification mechanisms.  The null equation
follows from the Gamma recurrence after shifting two Fourier-integral
contours by $\pm i$.  The shifted common Gamma factor has its nearest
poles at $\pm3i/2$; only the tanh multiplier contributes the two
simple poles at $\pm i/2$.  Their residue terms provide exactly the
subtracted $G_3^{(6)}(0)$, while the sinh multipliers contribute zero.
For the initial determinant, $\operatorname{Im}(x+i/2)^{2n}$ has
leading term $nx^{2n-1}$.  Continuous Cauchy--Binet therefore reduces
its sign to the product of the two alternants
\[
 \det(x_r^{2n-1})_{n,r=1}^3,
 \qquad \det(P_j(x_r))_{r,j=1}^3.
\]
After extracting $\prod_r x_r\tanh(\pi x_r)$, the first is a
Vandermonde in $x_r^2$ and the second is minus twice a Vandermonde
in $\cosh(\pi x_r)$.  Both variable maps are strictly increasing
on $(0,\infty)$; integration against $\prod_rW(x_r)$ is therefore
strictly negative.  This proves $\det T<0$ without a numerical
normalization constant.

\subsection{An explicit density bound}

Let
\begin{equation}\label{gs:scales}
 \lambda_j=j\pi,\qquad A=8(2\pi)^6,
 \qquad a_1=a_2=A/2,\quad a_3=A.
\end{equation}

\begin{lemma}[Uniform Gamma density estimate]\label{gs:density}
For every $x\ge2$ and $1\le j\le3$,
\begin{equation}\label{gs:density-bound}
 P_j(x)W(x)=a_jx^{-3}e^{-\lambda_jx}(1+\delta_j(x)),
 \qquad |\delta_j(x)|\le\frac6{x^2}.
\end{equation}
Moreover $a_j>1$, and on $0\le x\le2$,
\begin{equation}\label{gs:small-density}
                  0\le P_j(x)W(x)/a_j<2^{54}.
\end{equation}
\end{lemma}

\begin{proof}
We use the classical complex Stirling remainder in DLMF
\S5.11(ii).  For $z=1/4+ix/2$, $|\arg z|<\pi/2$, so retaining
$1/(12z)$ in Stirling's logarithmic expansion gives
\[
 \log\Gamma(z)=(z-\tfrac12)\log z-z+\tfrac12\log(2\pi)
                         +\frac1{12z}+\rho(z),
 \qquad |\rho(z)|\le\frac1{90|z|^3}.
\]
Indeed the first omitted term has magnitude $1/(360|z|^3)$,
and the secant factor is at most $4$.
Set $t=1/(2x)$ and
\[
 L(x)=\log\frac{W(x)}{Ax^{-3}e^{-3\pi x}}.
\]
Taking twelve times the real part and simplifying the elementary
logarithm yields
\[
 L(x)=-\frac32\log(1+t^2)
       +3\left(\frac{\arctan t}{t}-1\right)
       +\frac{4t^2}{1+t^2}+12\operatorname{Re}\rho(z).
\]
The inequalities $\log(1+t^2)\le t^2$ and
$0\le1-\arctan(t)/t\le t^2/3$ give
\[
 |L(x)|\le\frac{13}{8x^2}+\frac{16}{15x^3}
        \le\frac{259}{120x^2}<\frac9{4x^2}\qquad(x\ge2).
\]
Since $e^{9/16}<2$,
\[
 \left|\frac{W(x)}{Ax^{-3}e^{-3\pi x}}-1\right|
        \le\frac9{2x^2}.
\]
The elementary multipliers differ from their leading exponentials by
at most $2e^{-2\pi x}$.  More precisely their normalized factors are
$1-e^{-4\pi x}$, $1-e^{-2\pi x}$, and
$(1-e^{-2\pi x})/(1+e^{-2\pi x})$, respectively, and each lies
between zero and one.  Since $\pi>3$ and
$e^{6x}\ge(6x)^2/2$,
\[
                    2e^{-2\pi x}\le\frac1{9x^2}.
\]
Combining the two errors proves \eqref{gs:density-bound}, with
room to spare.

Euler's Gamma integral gives
$|\Gamma(1/4+ix/2)|\le\Gamma(1/4)<5$, because the contributions
from $(0,1)$ and $(1,\infty)$ are less than $4$ and $1$.
For $0\le x\le2$ all three multipliers are less than $e^{16}$.
Using $e<3$, $5^{12}<2^{28}$, $3^{16}<2^{26}$, and $a_j>1$
proves \eqref{gs:small-density}.
\end{proof}

\subsection{A completely explicit shifted Laplace estimate}

For an even integer $N\ge4096$ and $\lambda\in\{\pi,2\pi,3\pi\}$,
write
\[
 S_N(\lambda)=\frac{\Gamma(N-2)}{\lambda^{N-2}},\quad
 S_N^{(1)}(\lambda)=\frac{\Gamma(N-3)}{\lambda^{N-3}},\quad
 S_N^{(2)}(\lambda)=\frac{\Gamma(N-4)}{\lambda^{N-4}}.
\]

\begin{lemma}[Effective two-term phase estimate]\label{gs:Laplace}
For $1\le j\le3$ and every even $N\ge4096$,
\begin{equation}\label{gs:Laplace-bound}
 \left|\frac1{a_j}\int_0^\infty(x+i/2)^NP_j(x)W(x)\,dx
 -e^{i\lambda_j/2}
     \left(S_N(\lambda_j)+\frac{3i}{2}S_N^{(1)}(\lambda_j)\right)
 \right|\le16S_N^{(2)}(\lambda_j).
\end{equation}
\end{lemma}

\begin{proof}
The contribution of $[0,2]$ is at most $2^{55}3^N$ by
\eqref{gs:small-density}.  On $[2,\infty)$ substitute the model
$x^{-3}e^{-\lambda x}$ and use \eqref{gs:density-bound}.
The density error is at most
\[
 6\int_2^\infty|x+i/2|^Nx^{-5}e^{-\lambda x}\,dx.
\]
On $x\ge\sqrt N$,
\[
 |x+i/2|^N\le x^N\exp\left(\frac{N}{8x^2}\right)<2x^N,
\]
so this part is at most $12S_N^{(2)}(\lambda)$.  On
$2\le x\le\sqrt N$, its bound is
\[
                       \sqrt N(N+1)^{N/2}.
\]
Here we used $6x^{-5}\le6/32<1$ and the length at most $\sqrt N$.

For the model integral set $y=x+i/2$ and shift the ray to the
positive real axis beginning at $2$.  The integrand
\[
       e^{i\lambda/2}y^N(y-i/2)^{-3}e^{-\lambda y}
\]
has no pole in this strip.  The short vertical segment at $2$
is bounded by $3^N/16$, and the segment at infinity vanishes.
For real $y\ge2$, the binomial series gives
\[
 \left|(1-i/(2y))^{-3}-1-\frac{3i}{2y}\right|
                  \le\frac5{2y^2}.
\]
To check the constant, write $u=1/(2y)\le1/4$; then
$\sum_{k\ge2}\binom{k+2}{2}u^k\le10u^2$ since its quotient
by $u^2$ is increasing and is less than $10$ at $u=1/4$.
Thus this model error is at most $(5/2)S_N^{(2)}(\lambda)$.
Extending the first two retained integrals from $[2,\infty)$
to $(0,\infty)$ contributes less than $2^N$.

It remains to absorb the compact and lower-range contributions.
Their sum is bounded by
\[
 R_N=2^{56}3^N+\sqrt N(N+1)^{N/2}.
\]
Since $\lambda<10$ and $(N-5)!\ge((N-5)/3)^{N-5}$,
\[
 S_N^{(2)}(\lambda)
 \ge\frac{(N-5)^{N-5}}{3^{N-5}10^{N-4}}.
\]
For every $N\ge4096$,
\[
 \frac{30\sqrt{N+1}}{N-5}\le\frac12,
 \qquad \frac{90}{N-5}\le\frac1{32},
 \qquad N^6\le2^{N-2}.
\]
The first follows by squaring and expanding a quadratic; the second
is immediate; the third holds at $N=4096$ and persists by induction
because $(1+1/N)^6<2$.  Therefore
\[
 \frac{\sqrt N(N+1)^{N/2}}{S_N^{(2)}}
       \le N^6 2^{-N}\le\frac14,
 \qquad
 \frac{2^{56}3^N}{S_N^{(2)}}
       \le2^{56}N^5 2^{-5N}<\frac14.
\]
Consequently $R_N<S_N^{(2)}$.  Adding the bounds gives
$12+5/2+1<16$, proving \eqref{gs:Laplace-bound}.
\end{proof}

The retained second term is essential.  At $j=2$ the leading
phase $e^{i\lambda_j/2}=-1$ has zero imaginary part, so discarding
that second term would give no nonzero leading coefficient for the
middle column.

\subsection{A determinant comparison valid for every large index}

Taking imaginary parts of \eqref{gs:Laplace-bound} gives, with
$N=2n\ge4096$,
\begin{align}
 t_{n,1}&=\frac{4a_1}{N!}S_N(\lambda_1)(1+\epsilon_{n,1}),\notag\\
 t_{n,3}&=-\frac{4a_3}{N!}S_N(\lambda_3)(1+\epsilon_{n,3}),
 &|\epsilon_{n,j}|&\le\frac{1600}{(N-4)^2}\quad(j=1,3),
 \label{gs:odd-errors}\\
 t_{n,2}&=-\frac{6a_2}{N!}S_N^{(1)}(\lambda_2)(1+\epsilon_{n,2}),
 &|\epsilon_{n,2}|&\le\frac{320}{3(N-4)}.
 \label{gs:even-error}
\end{align}
The estimates use only $\lambda_j<10$ and the exact Gamma ratios.

Let $b\ge2048$ and $n=b+i$ for $1\le i\le3$.
After extracting positive row and column factors and the column
signs $(+,-,-)$, the coefficient determinant is the determinant of
an array $E=(E_{ij})$ satisfying
\begin{equation}\label{gs:relative-matrix}
 E_{ij}=j^{-2i}(1+\eta_{ij}),\qquad |\eta_{ij}|<\frac1{25}.
\end{equation}
For precision, extract the positive row factors
$4\Gamma(2b+2i-2)\pi^{-2i}/\Gamma(2b+2i+1)$ and positive column
factors $a_j\lambda_j^{-2b+2}$ for $j=1,3$.
For $j=2$, multiply this column factor additionally by
$3\lambda_2/(2(2b-1))$.
The middle-column row correction is
\[
 \rho_i=\frac{2b-1}{2b+2i-3},\qquad
 0<\rho_i\le1,\qquad1-\rho_i\le\frac4{2b-1}.
\]
Thus its relative error is bounded by
\[
 \frac4{2b-1}+\frac{320}{3(2b-2)}<\frac1{25}\qquad(b\ge2048),
\]
while the outer-column bound follows at once from
\eqref{gs:odd-errors}.  The product of the three extracted signs is
positive.

For the exact limiting array $V_{ij}=j^{-2i}$ one computes
\begin{equation}\label{gs:V}
 \det V=-\frac5{1944},\qquad
 \operatorname{perm}V=\frac{289}{23328}.
\end{equation}
The permanent here is the sum of the six unsigned permutation
products.  By expanding both determinants,
\[
 |\det E-\det V|
 \le\operatorname{perm}V\left\{(26/25)^3-1\right\}
 =\frac{289}{23328}\frac{1951}{15625}<\frac5{1944}.
\]
The exact positive margin is $373661/364500000$.
Therefore $\det E<0$, and hence
\begin{equation}\label{gs:large-positive}
                         f_b>0\qquad(b\ge2048)
\end{equation}
by \eqref{gs:Gamma-bridge} and $\det T<0$.
This is an effective all-index conclusion, rather than a numerical
comparison at finitely many large indices.

\subsection{Exact finite completion}

The remaining range is $4\le b\le2047$.
The following two elementary recurrences give a finite certificate
without computing huge rational Bernoulli numbers.  From
\[
 \frac{x}{2}\cot\frac{x}{2}=1-\frac12\sum_{n\ge1}\alpha_nx^{2n}
\]
and the differential equation of the cotangent, one obtains
\begin{equation}\label{gs:alpha-recurrence}
 \alpha_1=\frac16,\qquad
 \alpha_n=\frac1{2(2n+1)}
       \sum_{k=1}^{n-1}\alpha_k\alpha_{n-k}\quad(n\ge2).
\end{equation}
Combine this recurrence with \eqref{gs:recurrence} modulo
\[
                         \ell=2147483647=2^{31}-1.
\]
The supplied script proves primality by trial division through
$\lfloor\sqrt\ell\rfloor$.  Every denominator is invertible in
this field: for the needed indices $n\le2049$ the factors
$2(2n+1)$ are smaller than $\ell$, and
\[
 d_s=2\frac{(2s)!}{(2s-6)!}\frac{2s-3}{s}
\]
for $4\le s\le2050$ has no factor divisible by $\ell$.
Induction therefore identifies every computed residue with the
reduction of its exact rational coefficient.

\begin{proposition}[Finite rational nonvanishing certificate]
\label{gs:finite}
For each integer $4\le b\le2047$, the determinant $f_b$ has a
nonzero reduction modulo $2147483647$.  Hence $f_b\ne0$ over
$\mathbb Q$ throughout this finite interval.
\end{proposition}

\begin{proof}[Exact computational proof]
Run \texttt{verify\_gap\_six.py}.  It constructs all needed
$\alpha_n$ by \eqref{gs:alpha-recurrence}, constructs the three
fundamental columns by \eqref{gs:recurrence}, and expands each
three-by-three determinant by its six signed products.  All 2,044
residues are nonzero; the complete indexed list is in
\texttt{gap\_six\_certificate.json}.  Its canonical compact JSON
residue list has SHA-256
\begin{center}\small\ttfamily
90ec1b66f919cd00404dd92f7d1f9483ceb9\\
d53bfa5893735db62c797113fea6
\end{center}
The hash identifies the record, while the recurrences and the nonzero
residues constitute the mathematical certificate.

As independent normalization checks, all 2,049 cotangent coefficients
agree with formal division of cosine by the normalized sine series,
and the bridge \eqref{gs:bridge} agrees with direct exact determinants
of the original Bernoulli matrices for $4\le b\le20$.  The latter
comparisons check the omitted-index convention and every row/column
factor against \eqref{gs:entry}.  They are additional checks; the
finite-field recurrence itself proves the entire finite interval.
\end{proof}

Combining \eqref{gs:large-positive}, Proposition~\ref{gs:finite}, and
\eqref{gs:bridge} proves Theorem~\ref{gs:main}.  The general
unrestricted unequal-even-order conjecture remains open beyond the
families now covered.  In particular, proving all gaps of six does
not by itself prove the gap-eight family.

\subsection{An all-order family of rational polygamma-tail evaluations}

The rank theorem turns the source block calculus into an exact
ordinary convergent identity for every gap-six pair.  Let
$(\lambda_b,\ldots,\lambda_{2b-1})$ be the first row of $M_b^{-1}$,
and put
\begin{align}
 K_b(y)&=\sum_{j=b}^{2b-1}\lambda_jy^{2j+1},\notag\\
 P_b(x)&=K_b(x+\tfrac12)-K_b(x-\tfrac12)
                  =\sum_{r=0}^{q-1}w_rx^{2r}.\label{gs:polynomial}
\end{align}
All coefficients are explicitly specified rational numbers; the
inverse exists for every $b\ge4$ by Theorem~\ref{gs:main}.
For $h=0,1,2,3$ define
\[
 c_{v,0}=\frac1{v-1},\qquad
 c_{v,h}=\frac{B_{2h}(1/2)}{(2h)!}(v)_{2h-1}\quad(h\ge1),
 \qquad u_h=\sum_{e=0}^hc_{p,e}c_{q,h-e}.
\]
Here $(v)_r$ denotes the rising factorial.  Define the even sextic
subtraction polynomial
\begin{equation}\label{gs:subtraction}
 S_b(x)=\sum_{h=0}^3\left(\sum_{e=0}^hu_e\,w_{q-1-h+e}\right)
                                      x^{6-2h}.
\end{equation}

\begin{corollary}[Rational evaluations for every even gap of six]
\label{gs:ordinary}
For $b\ge4$, $p=2b-6$, $q=2b$, and $x_n=n+1/2$,
\begin{equation}\label{gs:ordinary-identity}
 \boxed{\displaystyle
 \sum_{n=0}^{\infty}\left\{
 \frac{P_b(x_n)}{(p-1)!(q-1)!}
       \psi^{(p-1)}(n+1)\psi^{(q-1)}(n+1)-S_b(x_n)
                 \right\}=-\frac12.}
\end{equation}
This sum converges absolutely.  Its polynomial $P_b$ has exact degree
$2q-2$, with leading coefficient
\begin{equation}\label{gs:leading-coefficient}
 w_{q-1}=\frac{4b(4b-1)f_{b-1}}{d_{b+3}f_b}\ne0,
\end{equation}
where the auxiliary initial determinant is
$f_3=-1/638788869734400000$.
\end{corollary}

\begin{proof}
The first inverse row gives
$\sum_{j=b}^{2b-1}\lambda_jC_j=Z_0=-1/2$.
The coefficients of the centered difference in
\eqref{gs:polynomial} are exactly
\[
 w_r=\sum_{j=\max\{b,r\}}^{2b-1}
       \lambda_j\binom{2j+1}{2r}2^{-2(j-r)},
\]
so \eqref{gs:backward} gives $\sum_rw_rD_r=-1/2$.
The centered tail expansion gives
\[
 A_p(n)A_q(n)=\sum_{h=0}^3u_hx_n^{2-p-q-2h}
                       +O(x_n^{-p-q-6}).
\]
Since $\deg P_b\le2q-2$ and $q-p=6$, its nonnegative powers
are exactly \eqref{gs:subtraction}; the remaining summand is
$O(x_n^{-2})$.  The continued sums of $1,x_n^2,x_n^4,x_n^6$
are zero because $\zeta(-2h,1/2)=0$.  Thus the ordinary
absolutely convergent sum equals the displayed moment combination.
Finally $A_p(n)=\psi^{(p-1)}(n+1)/(p-1)!$ for even $p$.

For \eqref{gs:leading-coefficient}, suppose first that $b\ge5$.
The cofactor defining $(M_b^{-1})_{0,b-1}$ deletes the final row
and the first column of $M_b$.  Reindexing $j'=j-1$, $k'=k-1$,
$p'=p-2$, and $q'=q-2$ in \eqref{gs:entry} shows that this
minor is exactly $M_{b-1}$: both binomial pairs, the Bernoulli
index, and the omitted-index convention are unchanged.  Hence
\[
 \lambda_{2b-1}=(-1)^{b-1}\frac{\det M_{b-1}}{\det M_b}
              =\frac{4bf_{b-1}}{d_{b+3}f_b}
\]
by \eqref{gs:bridge}.  The centered difference multiplies its
highest coefficient by $4b-1$, proving the formula and
nonvanishing for $b\ge5$.  At $b=4$ a direct exact calculation
gives $\lambda_7=84/83$, $d_7=6795360$,
$f_4=-83/22789201752651350016000000$, and the stated value of
$f_3$; these verify the same formula without invoking a
zero-order tail.  This completes the proof.
\end{proof}

\paragraph{Further research.}
The same Gamma basis supplies $m$ distinct exponential scales for
any fixed half-gap $m$.  Making all of its phase estimates effective
would give a numerical finite-completion threshold for each fixed
gap.  The next concrete target is $m=4$, where the two even modes
require their first nonzero phase corrections and the determinant
comparison involves a four-by-four reciprocal-square Vandermonde.
A different question is whether a direct recurrence invariant can
prove gap-six nonvanishing without the finite completion; the initial
sign transition shows why a fixed-sign assertion from the first
admissible index is false.
